\documentclass[12pt, a4paper, twoside, english, brazil]{article}
\usepackage[top=2cm, bottom=2cm, left=2cm, right=2cm]{geometry}

\usepackage{mathtools}
\usepackage{amsfonts}
\usepackage{amssymb}
\usepackage{amsmath}
\usepackage[utf8]{inputenc}
\usepackage[portuguese]{babel}
\usepackage[T1]{fontenc}

\usepackage{enumerate}
\usepackage{multicol}
\usepackage{tikz}
\usepackage{pgfplots}

\newenvironment{solution}[1][]{\noindent\textbf{Solução#1: \\} }

% -----------------------------------------------------------------------------

\title{Matemática Discreta e Lógica Matemática - Lista 01}
\author{Davi Reis - Funções de Primeiro Grau (10/04/2026)}
\date{}

\begin{document}

\maketitle

\begin{enumerate}[\bfseries 1.]
    % Questão 01
  \item Seja a função $p(s)=7+2.5s$, que calcula o preço (R\$), em
    função do espaço(s) percorrido(s) em $km$. Quanto se irá pagar ao
    andar 3 $km$ e 9 $km$.

    \begin{solution}\\
      Preço para 3: $p(3)=7+2.5\cdot 3=7+ 7.5 = 14.5$\\
      Preço para 9: $p(9)=7+2.5\cdot 9=7+ 22.5 = 29.5$
    \end{solution}

    % Questão 02
  \item Esboce o gráfico das funções $f(x)=2x+3$, $g(x)=4x+3$,
    $p(x)=-3x-2$ e $q(x)=-9x-2$.

    \begin{solution}
      \begin{tikzpicture}
        \begin{axis}[
            axis lines = center,
            xlabel = \(x\),ylabel ={\(y\)},
            xtick distance=0.5, ytick distance=1,
            xmax=1, xmin=-3,
            ymax=5, ymin=-5,
            grid,
            legend pos=north west
          ]

          \addplot [domain=-2:2,color=red]{2 * x + 3};
          \addplot [domain=-2:2,color=green]{4 * x + 3};
          \addplot [domain=-2:2,color=blue]{-3 * x - 2};
          \addplot [domain=-2:2,color=purple]{-9 * x - 2};

          \addlegendentry{$f(x)$}
          \addlegendentry{$g(x)$}
          \addlegendentry{$p(x)$}
          \addlegendentry{$q(x)$}
        \end{axis}
      \end{tikzpicture}
    \end{solution}

    % Questão 03
  \item Com base na Figura abaixo, encontre a reta $r(x)=ax+b$ que
    passa por esses 2 pontos.

    \begin{multicols}{2}
      \begin{tikzpicture}
        \begin{axis}[
            nodes near coords,
            axis lines = center,
            xlabel = \(x\),ylabel ={\(y\)},
            xmin=0, xmax=6,
            ymin=0, ymax=7,
            ytick distance=1,
            grid
          ]
          \addplot+ [
            only marks,
            point meta=explicit symbolic,
          ] coordinates {
            (3,6) [A]
            (5,2) [B]
          };
        \end{axis}

      \end{tikzpicture}

      \begin{solution}
        Dados os pontos podemos montar o sistema de duas variáveis:

        \noindent

        \begin{equation}
          \begin{cases}
            3a+b=6\\
            5a+b=2
          \end{cases}
        \end{equation}

        Subtraindo a primeira da segunda obtemos:\\
        $-2a=4\implies a=-2$\\
        e substituindo em na segunda, obtemos:\\
        $5\cdot (-2)+b=2\implies b=12$

        Portanto: $r(x)=-2x+12$
      \end{solution}
    \end{multicols}

    \newpage

    % Questão 04
  \item Encontre a raiz da função $h(x)=-2x+6$ e desenhe o gráfico se
    baseando no coeficiente linear e sua raiz.

    \begin{solution}
      Sabendo que a raiz de $h(x)$ é quando $h(x)=0$ temos: $0=-2x+6$,
      logo a raiz é $x=\frac{-6}{-2}=3$

      \begin{tikzpicture}
        \begin{axis}[
            axis lines = center,
            xlabel = \(x\),ylabel ={\(y\)},
            xmax=5, xmin=-2,
            ymax=8, ymin=-2,
            grid,
          ]
          \addplot [domain=-10:10,color=blue]{-2 * x + 6};
        \end{axis}
      \end{tikzpicture}
    \end{solution}

    % Questão 05
  \item Um objeto lançado verticalmente para cima, tem sua velocidade
    ($m/s$) descrita pela função (modelo) $v(t)=18t-216$.
    \begin{enumerate}
      \item Qual a sua velocidade inicial?
      \item Qual a velocidade em $t=6s$ e $t=14s$
      \item Quando o objeto alcança a altura máxima?
    \end{enumerate}

    \begin{solution}
      \hbox

      \begin{enumerate}
        \item Tomando $t=0$ obtemos $v_0=-216$, ou $216 m/s$ para cima.
        \item $v(6)=-108$ e $v(14)=36$, ou 108 $m/s$ para cima e 36 $m/s$ para
          baixo, respectivamente.
        \item Quando a velocidade é nula. Ou seja $v(t)=0\implies t=12 s$.
      \end{enumerate}
    \end{solution}

\end{enumerate}
\end{document}
